\documentclass[10pt]{article}
\author{Yi Gu}


\usepackage{mathrsfs}

%\usepackage{amssymb,amsthm,amsmath,amssymb,latexsym,amsfonts,txfonts, pxfonts,wasysym,stmaryrd}
\textheight=230mm \textwidth=140mm
\topmargin=-1.5cm
\oddsidemargin=1.2cm
\evensidemargin=1.2cm
\setlength\parindent{0pt}
\parskip=7pt
%\setcounter{page}{1}
%\setlength{\baselineskip}{24pt} \baselineskip=24pt
\usepackage{amsmath,amsthm,amssymb}
\usepackage{graphicx}
\usepackage{bbm}
\usepackage{xcolor}
\usepackage[utf8]{inputenc}
\usepackage[english]{babel}
\usepackage{enumitem}
\usepackage{hyperref}
\usepackage{color}
\usepackage{colortbl}
\usepackage{tikz}
%\usepackage{vector}

\hypersetup{
    colorlinks=true,
    linkcolor=blue,
    filecolor=magenta,      
    urlcolor=cyan,
    pdftitle={Sharelatex Example},
    citecolor=red,
    % bookmarks=true,
}


\font\tencmmib=cmmib10 \skewchar\tencmmib '60
\newfam\cmmibfam
\textfont\cmmibfam=\tencmmib
\def\Ph{\mathchar'4010}
\def\Ps{\mathchar'4011}
\font\tensmc=cmcsc10
\def\smc{\tensmc}
\def\srm{\sevenrm }
\font\tenmsb=msbm10
\font\eightmsb=msbm10 scaled 1100
\def\Bbb#1{\hbox{\tenmsb#1}}
\def\Bbbb#1{\hbox{\eightmsb#1}}
\def\bbox{\quad\hbox{\vrule \vbox{\hrule \vskip2pt \hbox{\hskip2pt
\vbox{\hsize=1pt}\hskip2pt} \vskip2pt\hrule}\vrule}}
\def\lessim{\ \lower4pt\hbox{$
\buildrel{\displaystyle <}\over\sim$}\ }
\def\gessim{\ \lower4pt\hbox{$\buildrel{\displaystyle >}
\over\sim$}\ }
\def\n{\noindent}
\def\P{{\cal P}}
\def\si{\bar{\sigma}}
\def\E{{\cal E}}
\def\LL{{\cal L}}
\def\FF{{\cal F}}
\def\SS{{\cal S}}
\def\C{{\cal C}}
\def\D{{\cal D}}
\def\H{{\cal H}}
\def\X{{\cal X}}
\def\Y{{\cal Y}}
\def\A{{\cal A}}
\def\B{{\cal B}}
\def\O{{\cal O}}
\def\M{{\cal M}}
\def\I{{\cal I}}
\def\eps{{\varepsilon}}
\def\ch{{\mbox{\rm ch}\hspace{0.4mm}}}
\def\sh{{\mbox{sh}}}
\def\th{{\mbox{th}}}
\def\Av{{\mbox{\rm{Av}}}}
\def\n{\noindent}
\def\bbe{\vskip .15pc\hangindent=.5pc\hangafter=1}
\def\goin{\to\infty}
\def\qed{\hfill\break\rightline{$\bbox$}}



\newcommand{\e}{\mathbb{E}}
\newcommand{\p}{\mathbb{P}}
\newcommand{\lra}{\leftrightarrow}
\newcommand{\nlra}{\nleftrightarrow}
\newcommand{\bet}{\begin{theorem}}
\newcommand{\ent}{\end{theorem}}
\newcommand{\Var}{\mathrm{Var}}
\newcommand{\g}{\Gamma}
\newcommand{\de}{\delta}
\newcommand{\norm}[1]{\left\lVert#1\right\rVert}
\newcommand{\gy}[1]{\textbf{\color{red}[Yi: #1]}}
\newcommand{\olsi}[1]{\,\overline{\!{#1}}}
\newcommand\y{\cellcolor{blue!20}}
\newcommand\x{\cellcolor{red!20}}
\DeclareMathOperator{\Tr}{Tr}

\newtheorem{lemma}{\bf Lemma}
\newtheorem{claim}{\bf Claim}
\newtheorem{definition}{\bf Definition}
\newtheorem{theorem}{\bf Theorem}
\newtheorem{corollary}{\bf Corollary}
\newtheorem{remark}{\bf Remark}
\newtheorem{example}{\bf Example}
\newtheorem{exercise}{\bf Exercise}
\newtheorem{proposition}{\bf Proposition}
\newtheorem{question}{\bf Question}
\newtheorem{acknowledgements}{\bf Acknowledgements}

\newenvironment{Proof of lemma}{\noindent{\bf Proof of Lemma}}{\hfill$\Box$\newline}
\newenvironment{Proof of claim}{\noindent{\bf Proof of claim}}{\hfill$\Box$\newline}
\newenvironment{Proof of theorem}{\noindent{\bf Proof of Theorem}}{\hfill\scriptsize{$\Box$}\newline}
\newenvironment{Proof of theorems}{\noindent{\bf Proof of Theorems}}{\hfill$\Box$\newline}
\newenvironment{Proof of proposition}{\noindent{\bf Proof of Proposition}}{\hfill$\Box$\newline}
\newenvironment{Proof of propositions}{\noindent{\bf Proof of Propositions}}{\hfill$\Box$\newline}
\newenvironment{Proof of exercise}{\noindent{\it Proof of Exercise:}}{\hfill$\Box$}
\newenvironment{Acknowledgements}{\noindent{\bf Acknowledgements}}



\begin{document}

Let $m$ be a fixed positive integer and $N$ a positive integer that will later go to $\infty$. Let $M$ be a symmetric $\binom{N}{m}$-by-$\binom{N}{m}$ matrix. The rows and columns are indexed by the subsets of $[N] = \{1, \cdots, N\}$ that have exactly $m$ elements. (Indeed there are obviously $\binom{N}{m}$ many of them.) 


Let $S$ and $T$ be two subsets of $[N]$ with $m$ elements. Finally, let $f_N(l): [m] \to \mathbb{N}$ be a function with $N$ as parameter. $f_N$ The matrix $M$ satisfies the following:

If $|S \cap T| = l$, then the entry at row $S$ and column $T$ has value $f_N(l)$.

We want to do asymptotic analysis on the following integral by letting $N \to \infty$:
\begin{align} \label{target}
    \binom{N}{m}^{-1}\log \int_{[0,\infty)^{\binom{N}{m}}} \exp \left(x^T M x\right)dx.
\end{align}

The setup may seem confusing at the first. To give some intuition, consider the special case that $m =1$. Then we have a normal $N$-by-$N$ matrix that can be written as
\begin{align} \label{matrix_sum}
    M = a_N\cdot Id + b_N \mathbbm{1}^T\mathbbm{1},
\end{align}
where $\mathbbm{1}$ is the vector with all entries 1. (\ref{target}) can be written as
\begin{align}
    \frac{1}{N}\log \int_{[0,\infty)^N} \exp \left(b_N \norm{x}_1^2 + a_N\norm{x}_2^2\right) dx.
\end{align}
The above expression can be written as the expectation of a function of $\norm{X}_1$, where $X = (|X_1|,\cdots,|X_N|)$ and $X_i$ are i.i.d standard Gaussian. We can use Laplace method and compute the asymptotics.

The problem becomes drastically more difficult when $m>1$. Take $m=2$ for instance, (\ref{matrix_sum}) becomes
\begin{align} 
    M = a_N\cdot Id + b_N \mathbbm{1}^T\mathbbm{1} + c_NM',
\end{align}
where $M'$ satisfies that the entry at row $\{a,b\}$ and column $\{c,d\}$ is 1 if $\{a,b\} \cap \{c,d\} = 1$ and 0 otherwise. ($m=2$ means the matrix is indexed by 2-subsets.) We have
\begin{align}
    x^TMx = a_N\norm{x}_2^2 + b_N \norm{x}_1^2 + c_Nx^TM'x.
\end{align}
Unless we know the rate function of $\langle X, M'X \rangle$ then the integral in (\ref{target}) can be handled by simply applying the contraction principle. However, it is hard to compute the rate function of $\langle X, M'X \rangle$, and we don't know if it is possible.

Questions:
1. For case $m=2$, is it possible to derive a LDP for $\langle X, M'X \rangle$? If so, can this method be extended to any fixed positive integer $m$?

2. If the setup is too cumbersome. Then for general symmetric, positive-definite matrix $M$, do we have a systematic approach to compute (\ref{target})?

\begin{remark}
    The matrix $M$ is actually the inverse of a covariance matrix. I managed to compute the inverse of the covariance, but what waits for me after a road block is another road block.
\end{remark}

\end{document}